\documentclass[12pt,a4paper]{article}

% ---- Encoding & Font ----
\usepackage{iftex}
\ifPDFTeX
  \usepackage[utf8]{inputenc}
  \usepackage[T1]{fontenc}
\else
  \usepackage{fontspec}  % XeTeX/LuaTeX (tectonic): Unicode nativo (·, —, ¿, ª, §)
\fi
\usepackage[spanish]{babel}
\usepackage{csquotes}

% ---- Page layout ----
\usepackage[top=2.5cm, bottom=2.5cm, left=2.5cm, right=2.5cm]{geometry}
\usepackage{setspace}
\onehalfspacing

% ---- Math ----
\usepackage{amsmath, amssymb, amsthm}
\usepackage{mathtools}
\usepackage{physics}
\usepackage{esint}  % \oiint

% ---- Graphics ----
\usepackage{graphicx}
\usepackage{tikz}
\usepackage{float}
\usepackage{caption}

% ---- Tables ----
\usepackage{array}
\usepackage{booktabs}
\usepackage{tabularx}
\usepackage{colortbl}

% ---- Colours ----
\usepackage{xcolor}
\definecolor{notebg}{HTML}{FFF3CD}
\definecolor{noteborder}{HTML}{FFC107}
\definecolor{boxred}{HTML}{F8D7DA}
\definecolor{boxredborder}{HTML}{F5C6CB}

% ---- Boxes ----
\usepackage{mdframed}
\usepackage{tcolorbox}
\tcbuselibrary{most}

\newtcolorbox{keybox}[1][]{
  colback=blue!5,
  colframe=blue!70!black,
  arc=4pt,
  boxrule=1pt,
  fonttitle=\bfseries,
  title={#1}
}

\newtcolorbox{warnbox}[1][]{
  colback=boxred,
  colframe=boxredborder,
  coltitle=black!75,
  arc=4pt,
  boxrule=1pt,
  fonttitle=\bfseries,
  title={#1}
}

\newtcolorbox{notebox}[1][]{
  colback=notebg,
  colframe=noteborder,
  coltitle=black!75,
  arc=4pt,
  boxrule=1pt,
  fonttitle=\bfseries,
  title={#1}
}

% ---- Hyperlinks & TOC ----
\usepackage[hidelinks]{hyperref}
\usepackage{bookmark}

% ---- Enumerate ----
\usepackage{enumitem}

% ---- Miscellaneous ----
\usepackage{icomma}
\usepackage{siunitx}
\usepackage{textcomp}
\usepackage[bottom]{footmisc}

% ---- Diagramas (gráficas y esquemas) ----
\usepackage{pgfplots}
\pgfplotsset{compat=1.18}
\usetikzlibrary{babel}  % babel español activa `>`; esto evita romper las flechas `->`
% courses/ElecMag2024/Clases/assets/tikzstyles.tex
% ---------------------------------------------------------------------------
% Estilos compartidos para los diagramas del curso Electromagnetismo (2024).
% Fuente única del "look" visual (colores, grosores, escalas) para que todas
% las figuras (TikZ / pgfplots / CircuiTikz) sean consistentes entre clases.
%
% Es una COPIA de courses/Fisica3-2015/Clases/assets/tikzstyles.tex: mismo
% dominio, misma paleta. Copia y no symlink para que cada curso sea
% autocontenido y la paleta de Física III pueda evolucionar aparte.
%
% Uso: la clase debe cargar antes `tikz` y `pgfplots` (y `circuitikz` si la
%      clase tiene circuitos), y luego % courses/ElecMag2024/Clases/assets/tikzstyles.tex
% ---------------------------------------------------------------------------
% Estilos compartidos para los diagramas del curso Electromagnetismo (2024).
% Fuente única del "look" visual (colores, grosores, escalas) para que todas
% las figuras (TikZ / pgfplots / CircuiTikz) sean consistentes entre clases.
%
% Es una COPIA de courses/Fisica3-2015/Clases/assets/tikzstyles.tex: mismo
% dominio, misma paleta. Copia y no symlink para que cada curso sea
% autocontenido y la paleta de Física III pueda evolucionar aparte.
%
% Uso: la clase debe cargar antes `tikz` y `pgfplots` (y `circuitikz` si la
%      clase tiene circuitos), y luego % courses/ElecMag2024/Clases/assets/tikzstyles.tex
% ---------------------------------------------------------------------------
% Estilos compartidos para los diagramas del curso Electromagnetismo (2024).
% Fuente única del "look" visual (colores, grosores, escalas) para que todas
% las figuras (TikZ / pgfplots / CircuiTikz) sean consistentes entre clases.
%
% Es una COPIA de courses/Fisica3-2015/Clases/assets/tikzstyles.tex: mismo
% dominio, misma paleta. Copia y no symlink para que cada curso sea
% autocontenido y la paleta de Física III pueda evolucionar aparte.
%
% Uso: la clase debe cargar antes `tikz` y `pgfplots` (y `circuitikz` si la
%      clase tiene circuitos), y luego \input{../assets/tikzstyles.tex}
% (compilar desde el directorio de la clase para que la ruta relativa resuelva).
% ---------------------------------------------------------------------------

% ---- Paleta (alineada con las cajas del preámbulo: keybox/notebox/warnbox) ----
\definecolor{figblue}{HTML}{1F4E79}   % azul principal (líneas, keybox)
\definecolor{figred}{HTML}{C0392B}    % rojo de énfasis (warnbox)
\definecolor{figamber}{HTML}{B8860B}  % ámbar (notebox)
\definecolor{figgray}{HTML}{555555}   % auxiliares, ejes, guías

% ---- CircuiTikz: escala y grosor de los componentes ----
% Guarda \ifdefined: la electrostática (Clases 1-14) no carga `circuitikz`, y
% sin la guarda el \input revienta con `Undefined control sequence \ctikzset`.
% Cuando el curso llegue a circuitos, el bloque se activa solo.
\ifdefined\ctikzset
  \ctikzset{
    bipoles/thickness=1,
    resistors/scale=0.9,
    capacitors/scale=0.9,
  }
\fi

% ---- TikZ: estilos reutilizables ----
% Nota: las puntas se declaran como `-{Latex[...]}` y no como `->`. La forma
% con llaves NO contiene un `>` literal, así que es inmune al `>` activo que
% introduce babel español (ver CLAUDE.md §7.3.3).
\usetikzlibrary{arrows.meta,calc,angles,quotes}

\tikzset{
  figline/.style={figblue, line width=0.9pt},
  figaux/.style={figgray, dashed, line width=0.6pt},
  % vectores
  figvec/.style={figblue, line width=0.9pt, -{Latex[length=2.2mm,width=1.6mm]}},
  figvecres/.style={figred, line width=1.3pt, -{Latex[length=2.6mm,width=1.9mm]}},
  figvecaux/.style={figgray, line width=0.7pt, -{Latex[length=1.8mm,width=1.3mm]}},
  figaxis/.style={figgray, line width=0.6pt, -{Latex[length=2mm,width=1.4mm]}},
  % cargas puntuales
  figcarga/.style={circle, inner sep=0pt, minimum size=5.4pt, draw=none},
  figpos/.style={figcarga, fill=figred},
  figneg/.style={figcarga, fill=figblue},
  figneutra/.style={figcarga, fill=figgray},
  % etiquetas
  figlbl/.style={font=\small},
  figlblaux/.style={font=\footnotesize, figgray},
}

% ---- pgfplots: estilo base de las gráficas ----
\pgfplotsset{
  emfig/.style={
    width=11cm, height=6.5cm,
    axis lines=left,
    xlabel style={font=\small},
    ylabel style={font=\small},
    tick label style={font=\footnotesize},
    every axis plot/.append style={line width=1pt},
    grid=major,
    grid style={figgray!20},
    legend cell align=left,
    legend style={font=\footnotesize, draw=figgray!40},
  },
}

% (compilar desde el directorio de la clase para que la ruta relativa resuelva).
% ---------------------------------------------------------------------------

% ---- Paleta (alineada con las cajas del preámbulo: keybox/notebox/warnbox) ----
\definecolor{figblue}{HTML}{1F4E79}   % azul principal (líneas, keybox)
\definecolor{figred}{HTML}{C0392B}    % rojo de énfasis (warnbox)
\definecolor{figamber}{HTML}{B8860B}  % ámbar (notebox)
\definecolor{figgray}{HTML}{555555}   % auxiliares, ejes, guías

% ---- CircuiTikz: escala y grosor de los componentes ----
% Guarda \ifdefined: la electrostática (Clases 1-14) no carga `circuitikz`, y
% sin la guarda el \input revienta con `Undefined control sequence \ctikzset`.
% Cuando el curso llegue a circuitos, el bloque se activa solo.
\ifdefined\ctikzset
  \ctikzset{
    bipoles/thickness=1,
    resistors/scale=0.9,
    capacitors/scale=0.9,
  }
\fi

% ---- TikZ: estilos reutilizables ----
% Nota: las puntas se declaran como `-{Latex[...]}` y no como `->`. La forma
% con llaves NO contiene un `>` literal, así que es inmune al `>` activo que
% introduce babel español (ver CLAUDE.md §7.3.3).
\usetikzlibrary{arrows.meta,calc,angles,quotes}

\tikzset{
  figline/.style={figblue, line width=0.9pt},
  figaux/.style={figgray, dashed, line width=0.6pt},
  % vectores
  figvec/.style={figblue, line width=0.9pt, -{Latex[length=2.2mm,width=1.6mm]}},
  figvecres/.style={figred, line width=1.3pt, -{Latex[length=2.6mm,width=1.9mm]}},
  figvecaux/.style={figgray, line width=0.7pt, -{Latex[length=1.8mm,width=1.3mm]}},
  figaxis/.style={figgray, line width=0.6pt, -{Latex[length=2mm,width=1.4mm]}},
  % cargas puntuales
  figcarga/.style={circle, inner sep=0pt, minimum size=5.4pt, draw=none},
  figpos/.style={figcarga, fill=figred},
  figneg/.style={figcarga, fill=figblue},
  figneutra/.style={figcarga, fill=figgray},
  % etiquetas
  figlbl/.style={font=\small},
  figlblaux/.style={font=\footnotesize, figgray},
}

% ---- pgfplots: estilo base de las gráficas ----
\pgfplotsset{
  emfig/.style={
    width=11cm, height=6.5cm,
    axis lines=left,
    xlabel style={font=\small},
    ylabel style={font=\small},
    tick label style={font=\footnotesize},
    every axis plot/.append style={line width=1pt},
    grid=major,
    grid style={figgray!20},
    legend cell align=left,
    legend style={font=\footnotesize, draw=figgray!40},
  },
}

% (compilar desde el directorio de la clase para que la ruta relativa resuelva).
% ---------------------------------------------------------------------------

% ---- Paleta (alineada con las cajas del preámbulo: keybox/notebox/warnbox) ----
\definecolor{figblue}{HTML}{1F4E79}   % azul principal (líneas, keybox)
\definecolor{figred}{HTML}{C0392B}    % rojo de énfasis (warnbox)
\definecolor{figamber}{HTML}{B8860B}  % ámbar (notebox)
\definecolor{figgray}{HTML}{555555}   % auxiliares, ejes, guías

% ---- CircuiTikz: escala y grosor de los componentes ----
% Guarda \ifdefined: la electrostática (Clases 1-14) no carga `circuitikz`, y
% sin la guarda el \input revienta con `Undefined control sequence \ctikzset`.
% Cuando el curso llegue a circuitos, el bloque se activa solo.
\ifdefined\ctikzset
  \ctikzset{
    bipoles/thickness=1,
    resistors/scale=0.9,
    capacitors/scale=0.9,
  }
\fi

% ---- TikZ: estilos reutilizables ----
% Nota: las puntas se declaran como `-{Latex[...]}` y no como `->`. La forma
% con llaves NO contiene un `>` literal, así que es inmune al `>` activo que
% introduce babel español (ver CLAUDE.md §7.3.3).
\usetikzlibrary{arrows.meta,calc,angles,quotes}

\tikzset{
  figline/.style={figblue, line width=0.9pt},
  figaux/.style={figgray, dashed, line width=0.6pt},
  % vectores
  figvec/.style={figblue, line width=0.9pt, -{Latex[length=2.2mm,width=1.6mm]}},
  figvecres/.style={figred, line width=1.3pt, -{Latex[length=2.6mm,width=1.9mm]}},
  figvecaux/.style={figgray, line width=0.7pt, -{Latex[length=1.8mm,width=1.3mm]}},
  figaxis/.style={figgray, line width=0.6pt, -{Latex[length=2mm,width=1.4mm]}},
  % cargas puntuales
  figcarga/.style={circle, inner sep=0pt, minimum size=5.4pt, draw=none},
  figpos/.style={figcarga, fill=figred},
  figneg/.style={figcarga, fill=figblue},
  figneutra/.style={figcarga, fill=figgray},
  % etiquetas
  figlbl/.style={font=\small},
  figlblaux/.style={font=\footnotesize, figgray},
}

% ---- pgfplots: estilo base de las gráficas ----
\pgfplotsset{
  emfig/.style={
    width=11cm, height=6.5cm,
    axis lines=left,
    xlabel style={font=\small},
    ylabel style={font=\small},
    tick label style={font=\footnotesize},
    every axis plot/.append style={line width=1pt},
    grid=major,
    grid style={figgray!20},
    legend cell align=left,
    legend style={font=\footnotesize, draw=figgray!40},
  },
}


% ---- Title ----
\title{\textbf{Clase 22: Circuitos magnéticos} \\[0.3em]
        \large{Potencial escalar magnético, reluctancia e imanes permanentes}}
\author{Transcripto y expandido --- Curso Electromagnetismo (OpenFING)}
\date{}

\begin{document}

\maketitle
\thispagestyle{empty}
\newpage

\tableofcontents
\newpage

\section{Problemas magnéticos con condiciones de borde}

El ferromagnetismo, con sus curvas de histéresis y relaciones no lineales, da
en general problemas difíciles. Esta clase muestra situaciones en que se puede
avanzar con aproximaciones razonables, algunas con aplicaciones tecnológicas
importantes. La primera es parecida a la electrostática: hallar el campo en
regiones con condiciones en la frontera.

\subsection{Potencial escalar magnético}

Para corrientes estacionarias las ecuaciones son

\[
\div\vec B = 0, \qquad \curl\vec H = \vec J,
\]
con $\vec J$ la densidad de corriente \textbf{de conducción}, más la ecuación
constitutiva del o de los materiales, que relaciona $\vec B$ con $\vec H$.

Muchas veces la corriente está concentrada en algunas regiones. En las regiones
\textbf{sin corriente de conducción}, $\vec J = 0$ y entonces $\curl\vec H = 0$:
$\vec H$ deriva de un potencial,

\begin{keybox}
\[
\boxed{\ \vec H = -\grad\phi^*\ }
\]
\end{keybox}
el \textbf{potencial escalar magnético}. El asterisco es la notación del libro, para
no confundirlo con el potencial electrostático.

\begin{notebox}
A diferencia del campo eléctrico, que en electrostática \textbf{siempre} deriva de
un potencial, ésta es una situación particular, no genérica: vale sólo donde
no hay corriente de conducción. Corrientes de magnetización sí puede haber.
\end{notebox}

Sin ecuación constitutiva no se puede seguir; se ven dos casos.

\subsection{Caso 1: material lineal}

Exactamente lineal no es casi ninguno interesante, pero hay materiales cuya
curva de histéresis es muy angosta ---el \textbf{hierro dulce}, por ejemplo--- y para
los que $\vec B = \mu\vec H$ es una aproximación aceptable («más o menos»,
no tan buena como en el caso eléctrico). Con $\mu$ constante en cada material,
lejos de los bordes y donde no hay corriente:

\[
0 = \div\vec B = \mu\,\div\vec H = -\mu\,\laplacian\phi^*
\quad\Longrightarrow\quad
\boxed{\ \laplacian\phi^* = 0\ }
\]

Todos los métodos de la ecuación de Laplace se importan a los problemas
magnéticos. En particular vale en el vacío donde no hay corriente.

\begin{notebox}
Una diferencia práctica: en electrostática las cargas solían estar en
superficies; acá es común que la corriente esté distribuida en volumen. Laplace
vale sólo afuera de esas regiones.
\end{notebox}

\subsection{Caso 2: magnetización uniforme}

Un caso \textbf{no lineal} pero manejable: la magnetización dentro del material es,
en buena aproximación, \textbf{uniforme}. La relación general es
$\vec B = \mu_0(\vec H + \vec M)$, y si $\vec M$ es uniforme su divergencia es
nula:

\[
0 = \div\vec B = \mu_0\left(\div\vec H + \div\vec M\right)
= -\mu_0\,\laplacian\phi^* .
\]

Otra vez Laplace. Que no sea lineal se ve en que $\vec M$ no es proporcional a
$\vec B$: no depende de él, está fija.

\begin{center}
\renewcommand{\arraystretch}{1.35}
\small
\begin{tabularx}{\textwidth}{l>{\raggedright\arraybackslash}X>{\raggedright\arraybackslash}X}
\toprule
\textbf{Caso} & \textbf{Hipótesis} & \textbf{Resultado} \\
\midrule
1 & Material lineal, $\vec J = 0$ & $\laplacian\phi^* = 0$ \\
2 & $\vec M$ uniforme, $\vec J = 0$ & $\laplacian\phi^* = 0$ \\
\bottomrule
\end{tabularx}
\end{center}

\begin{warnbox}
\textbf{Ojo con las condiciones de borde}: no son las de electrostática, sino las
de la clase pasada ($B_n$ continua, $H_t$ continua salvo corriente
superficial de conducción).
\end{warnbox}

\subsection{Dos ejemplos}

Quedan como ejercicio (están en el práctico):

\begin{enumerate}
  \item \textbf{Esfera de permeabilidad $\mu$ en un campo uniforme $\vec B_0$}, rodeada
  de aire. Sin corriente de conducción en ningún lado. Es muy parecido a la
  esfera dieléctrica: simetría azimutal, desarrollos en polinomios de
  Legendre adentro y afuera, y las condiciones de borde.
  \item \textbf{Imán esférico con $\vec M$ uniforme}, sin campo externo aplicado. Lo que
  genera el campo es la magnetización espontánea, que es dato. Afuera hay aire
  (caso lineal con $\mu = \mu_0$) y adentro magnetización uniforme (caso 2).
  Con el eje $z$ según $\vec M$, el problema tiene simetría azimutal. Adentro
  las líneas son paralelas y equiespaciadas; afuera, curvas que se cierran.
\end{enumerate}

Con cilindros muy largos magnetizados a lo largo del eje se puede hacer lo
mismo en el plano transversal, con los métodos ya vistos.

\section{Circuitos magnéticos}

Otra situación, muy importante en la práctica y en la tecnología.

\subsection{Núcleos de alta permeabilidad}

Un \textbf{núcleo de hierro} que se cierra sobre sí mismo ---un toroide, aunque podría
ser rectangular o de otra forma--- al que se le conectan cosas: \textbf{bobinados}
(uno o varios), tramos de \textbf{distintos materiales} A, B, C, o un pedacito de
aire, el \textbf{entrehierro}, que tiene que ser chico. Puede haber varias mallas.

Los materiales del núcleo, lineales o no, tienen $\mu\gg\mu_0$, y afuera hay
aire.

\begin{figure}[H]
  \centering
  % Clase 22 — un circuito magnético: núcleo cerrado de alta permeabilidad con un
% bobinado de N vueltas, tramos de materiales distintos y un entrehierro chico.
% Las líneas de B circulan por el núcleo; la curva C va por el medio del núcleo.
\begin{tikzpicture}[scale=1]

  % nucleo (anillo rectangular)
  \fill[figgray!30, even odd rule]
    (0,0) rectangle (6,4)  (1,1) rectangle (5,3);
  \draw[figline] (0,0) rectangle (6,4);
  \draw[figline] (1,1) rectangle (5,3);
  % tramo de otro material (B) arriba
  \fill[figamber!30] (2.4,3) rectangle (3.6,4);
  \draw[figamber, line width=0.8pt] (2.4,3) -- (2.4,4) (3.6,3) -- (3.6,4);
  \node[figlbl, fill=figamber!30, inner sep=1pt] at (3.0,3.72) {B};
  \node[figlbl] at (5.78,2.85) {A};
  % entrehierro en el lado inferior
  \fill[white] (2.85,-0.05) rectangle (3.15,1.05);
  \draw[figline] (2.85,0) -- (2.85,1) (3.15,0) -- (3.15,1);
  \node[figlblaux, anchor=north] at (3.0,-0.1) {entrehierro};

  % bobinado en el lado izquierdo
  \foreach \yy in {1.2,1.5,1.8,2.1,2.4,2.7} {
    \draw[figred, line width=1pt] (-0.12,\yy) -- (1.12,\yy+0.12);
  }
  \draw[figred, line width=1pt] (-0.12,1.2) -- (-0.7,1.2);
  \draw[figred, line width=1pt] (-0.12,2.82) -- (-0.7,2.82);
  \node[figlbl, text=figred, anchor=east] at (-0.75,2.0) {$N,\ I$};

  % curva C por el medio
  \draw[figblue, dashed, line width=0.8pt] (0.5,0.5) rectangle (5.5,3.5);
  \draw[figvec] (5.5,1.6) -- (5.5,2.3);
  \node[figlbl, text=figblue] at (4.3,0.25) {$C$};

\end{tikzpicture}

  \caption{Un circuito magnético: núcleo cerrado de alta permeabilidad, un bobinado de $N$ vueltas, un tramo de otro material y un entrehierro chico. La curva $C$ de la ley de las mallas va por el medio del núcleo.}
\end{figure}

\textbf{La analogía clave.} En un cable con aire alrededor, los electrones no se
escapan del cable: la conductividad del cobre es muchísimo mayor que la del
aire. Acá pasa lo mismo con el campo magnético: si $\mu\gg\mu_0$, las líneas que
genera un bobinado \textbf{circulan dentro del núcleo} y no se van al aire.

¿Cuánto se escapa? Con $\mu = 4000\,\mu_0$ ---valor totalmente verosímil para
muchos ferromagnéticos---, aproximadamente una línea de cada 4000. Donde la
aproximación empeora es en el \textbf{entrehierro}, donde hay $\mu_0$; si es muy
chico, el campo «no tiene tiempo» de irse por los costados.

\subsection{Las aproximaciones y su calidad}

\begin{center}
\renewcommand{\arraystretch}{1.35}
\small
\begin{tabularx}{\textwidth}{l>{\raggedright\arraybackslash}X>{\raggedright\arraybackslash}X}
\toprule
\textbf{Aproximación} & \textbf{Qué supone} & \textbf{Calidad} \\
\midrule
De circuito & El campo circula por los núcleos de alto $\mu$ y no se dispersa en el aire & Muy buena (error $\sim\mu_0/\mu$), salvo en el entrehierro \\
Filiforme & $\vec B$ uniforme en cada sección, paralelo al núcleo y ortogonal a la sección & Regular: 5 o 10 \% de error \\
\bottomrule
\end{tabularx}
\end{center}

En los circuitos eléctricos las dos son casi exactas: la diferencia de
conductividad entre metal y aire es de unos veinte órdenes de magnitud, «blanco
o negro», y los cables son realmente finos. Acá la primera es buena pero no tan
buena, y la segunda es mala: un núcleo tiene centímetros de grosor ---no puede ser
muy fino, porque hay que enrollarle el bobinado--- y en las esquinas de un núcleo
rectangular el campo da vuelta y no es uniforme.

\begin{notebox}
«El profesor de electro se puso a hacer aproximaciones truchas. Y sí, pero a
veces la vida puede más»: es la razón del artillero, se tira con los obuses
que se tienen. Estas aproximaciones se usan muchísimo, pero hay que usarlas
con espíritu crítico y saber que, para algo preciso, hay que mejorarlas.
\end{notebox}

\subsection{Ley de los nudos}

El flujo magnético por una superficie cerrada es cero. Si en un nudo entra un
flujo $\Phi_1$ y se reparte en $\Phi_2$ y $\Phi_3$:

\begin{keybox}
\[
\boxed{\ \Phi_1 = \Phi_2 + \Phi_3\ }
\]
\end{keybox}

Es una muy buena aproximación: lo exacto es que el flujo total por la
superficie cerrada es nulo, y lo único que se desprecia es el flujo lateral que
se va al aire (1 en 4000), salvo en los entrehierros, donde la pérdida es
proporcional a la superficie lateral.

\subsection{Ley de las mallas y reluctancia}

Se toma una curva cerrada $C$ que recorre una malla por dentro del núcleo, y un
bobinado de $N$ vueltas con corriente $I$. Por Ampère en el material:

\[
\oint_C \vec H\cdot d\vec l = N\,I .
\]

El signo de cada término sale de la regla de la mano derecha según el sentido
en que se recorre $C$; con varios bobinados hay un término $N_k I_k$ por cada
uno.

Ahora sí las aproximaciones filiformes:

\begin{enumerate}
  \item $\vec H$ paralelo a $d\vec l$ en todo punto;
  \item $\vec B = \mu\vec H$ ---sin exigir linealidad: $\mu$ puede depender del campo---
  en un material homogéneo e isotrópico;
  \item $\vec B$ uniforme en la sección $A$, de modo que $\Phi = B\,A$.
\end{enumerate}

\[
\oint_C H\,dl = \oint_C \frac{B}{\mu}\,dl = \oint_C \frac{\Phi}{\mu A}\,dl .
\]

En una malla sin nudos el flujo es el mismo en todo el recorrido y sale de la
integral; $\mu$ y $A$ no, porque pueden cambiar a lo largo del circuito
(distintos materiales, tramos más gruesos o más finos):

\begin{keybox}
\[
\boxed{\ \Phi\,\mathcal{R} = N\,I,\qquad
\mathcal{R} = \oint_C \frac{dl}{\mu\,A}\ }
\]
\end{keybox}

$\mathcal{R}$ es la \textbf{reluctancia}. Con varias mallas hay una ecuación por
malla y, en cada tramo, el flujo que circula por él.

\begin{notebox}
Para disminuir los errores, la curva se toma por el medio del núcleo, que
promedia las diferencias de campo en la sección.
\end{notebox}

\subsection{Analogía con los circuitos eléctricos}

\begin{center}
\renewcommand{\arraystretch}{1.35}
\small
\begin{tabularx}{\textwidth}{>{\raggedright\arraybackslash}X>{\raggedright\arraybackslash}X}
\toprule
\textbf{Circuito eléctrico} & \textbf{Circuito magnético} \\
\midrule
Fuente $\mathcal{B}$ (batería) & $N I$ (bobinado) \\
Corriente $I$ & Flujo $\Phi$ \\
Densidad de corriente $\vec J$ & Campo $\vec B$ \\
Resistencia $R$ & Reluctancia $\mathcal{R}$ \\
$\mathcal{B} = R I$ & $N I = \mathcal{R}\,\Phi$ \\
Ley de los nudos & $\sum\Phi = 0$ en un nudo \\
\bottomrule
\end{tabularx}
\end{center}

Como con las resistencias, las reluctancias en serie se suman y en paralelo se
suman los inversos. Hasta el nombre se parece.

\subsection{La gran diferencia: no linealidad y resolución iterativa}

La pregunta del docente a la clase: aun con todas estas aproximaciones, ¿qué es
muy distinto del caso eléctrico? La resistencia \textbf{no depende} de la corriente:
la ley de Ohm es lineal. La reluctancia, en cambio, depende del flujo, porque
$\mu$ depende del campo en los materiales interesantes:

\[
\mathcal{R} = \mathcal{R}(\Phi) .
\]

Las ecuaciones se parecen pero el sistema es \textbf{no lineal}, y no se resuelve
con eliminación de Gauss. Se resuelve de manera \textbf{iterativa}:

\begin{enumerate}
  \item Se estima el rango de campo y se toma el $\mu$ correspondiente, constante.
  \item Se resuelve el problema, ya lineal, y se obtiene un campo.
  \item Se va a la tabla (o a la curva) $\mu(B)$ y se toma el $\mu$ de ese campo.
  \item Se vuelve a resolver, hasta que el campo supuesto y el calculado coincidan.
\end{enumerate}

No hace falta precisión de 1 por 1000: con las aproximaciones que se hicieron,
si $\mu$ cambia menos del 5 \% entre iteraciones, alcanza. Típicamente bastan un
par de iteraciones.

En electrotecnia se usan tablas reales; en el curso las curvas de histéresis se
reemplazan por tramos de recta.

\begin{center}
\renewcommand{\arraystretch}{1.35}
\small
\begin{tabularx}{\textwidth}{l>{\raggedright\arraybackslash}X>{\raggedright\arraybackslash}X}
\toprule
 & \textbf{Circuitos eléctricos} & \textbf{Circuitos magnéticos} \\
\midrule
Aproximaciones & Casi exactas & Similares, en general de peor calidad \\
Ecuaciones & Mallas y nudos & Similares en apariencia \\
Coeficientes & $R$ no depende de $I$ & $\mathcal{R}$ suele depender de $\Phi$ \\
Resolución & Sistema lineal & No lineal: iterativa \\
\bottomrule
\end{tabularx}
\end{center}

\section{Circuitos magnéticos con imanes permanentes}

En lugar de un bobinado, la fuente es un tramo de \textbf{imán}: un material con una
relación no trivial entre $\vec H$ y $\vec B$, típicamente con $\vec B\neq0$
cuando $\vec H\to0$ (magnetización espontánea). Hay que saber su curva.

\subsection{El circuito}

Un circuito de una malla con sección uniforme $A$: un tramo de imán entre los
puntos A y B, hierro dulce de B a C, un entrehierro de C a D y hierro dulce de
D a A. En el hierro dulce se supone que no hay histéresis y que
$\vec B = \mu\vec H$ es lineal; en el imán, no.

\begin{figure}[H]
  \centering
  % Clase 22 — circuito magnético con imán permanente, sin bobinados. Imán entre
% A y B, hierro dulce de B a C y de D a A, entrehierro entre C y D. Sección A
% uniforme. Domina la reluctancia del entrehierro.
\begin{tikzpicture}[scale=1]

  \fill[figgray!30, even odd rule] (0,0) rectangle (6,3.6)  (1,1) rectangle (5,2.6);
  \draw[figline] (0,0) rectangle (6,3.6);
  \draw[figline] (1,1) rectangle (5,2.6);

  % iman en el lado izquierdo
  \fill[figred!25] (0,1.0) rectangle (1,2.6);
  \draw[figred, line width=0.9pt] (0,1.0) -- (1,1.0) (0,2.6) -- (1,2.6);
  \node[figlbl] at (0.5,1.8) {imán};
  % entrehierro en el lado derecho
  \fill[white] (4.95,1.65) rectangle (6.05,1.95);
  \draw[figline] (5,1.65) -- (6,1.65) (5,1.95) -- (6,1.95);

  % puntos
  \fill (0.5,1.0) circle (1.6pt); \node[figlbl, anchor=north] at (0.5,0.95) {A};
  \fill (0.5,2.6) circle (1.6pt); \node[figlbl, anchor=south] at (0.5,2.65) {B};
  \fill (5.5,1.95) circle (1.6pt); \node[figlbl, anchor=west] at (6.1,2.1) {C};
  \fill (5.5,1.65) circle (1.6pt); \node[figlbl, anchor=west] at (6.1,1.5) {D};

  \node[figlblaux] at (3,3.85) {hierro dulce};
  \node[figlblaux] at (3,-0.25) {hierro dulce};
  \node[figlblaux, anchor=west] at (6.65,1.8) {entrehierro};

\end{tikzpicture}

  \caption{Circuito con imán permanente entre A y B, hierro dulce de B a C y de D a A, y entrehierro entre C y D, con sección uniforme.}
\end{figure}

\textbf{No hay bobinados}, así que la circulación de $\vec H$ en una curva cerrada
---por el medio del núcleo--- es \textbf{cero}:

\[
\int_A^B H\,dl + \int_B^C H\,dl + \int_C^D H\,dl + \int_D^A H\,dl = 0 .
\]

Con una sola malla el flujo $\Phi$ es el mismo en todo el circuito. En cada
tramo que no es el imán, $\int H\,dl = \Phi\,\mathcal{R}_{\text{tramo}}$:

\[
\int_A^B H\,dl + \Phi\left(\mathcal{R}_{BC} + \mathcal{R}_{CD} + \mathcal{R}_{DA}\right) = 0 .
\]

\subsection{Domina la reluctancia del entrehierro}

¿Qué reluctancia es mayor? La del \textbf{entrehierro}. Las reluctancias son
inversamente proporcionales a $\mu$: en el hierro dulce, $\mu\gg\mu_0$ y su
reluctancia es casi nula. El entrehierro es corto, pero típicamente no 4000
veces más corto que el hierro sino unas 20 veces; sólo con un entrehierro
miles de veces más chico que los tramos de hierro se compensaría.

\begin{notebox}[Por qué es una buena noticia]
Las aproximaciones sobre el hierro dulce
(lineal, filiforme) no son buenas, pero afectan a términos que casi no
cuentan. La del entrehierro, en cambio, se calcula muy bien:
\[
\mathcal{R}_{CD} = \frac{L_{CD}}{\mu_0\,A} .
\]
Es como en un circuito eléctrico: los tramos de hierro son los cables y el
entrehierro es el resistor.
\end{notebox}

El uso típico de un circuito así es \textbf{transportar} el campo del imán hasta la
zona del entrehierro, donde se lo quiere usar, sin que se disperse.

\subsection{Recta de desmagnetización y solución gráfica}

Despreciando $\mathcal{R}_{BC}$ y $\mathcal{R}_{DA}$ frente a $\mathcal{R}_{CD}$,
y suponiendo $H$ constante en el imán (de largo $L_{AB}$):

\[
H_m\,L_{AB} + \frac{\Phi\,L_{CD}}{\mu_0 A} = 0 .
\]

Como las secciones son iguales, el campo en el imán es el mismo que en el
entrehierro, $B_m = \Phi/A$ (si no, lo que sería igual son los flujos):

\begin{keybox}
\[
\boxed{\ B_m = -\,\mu_0\,\frac{L_{AB}}{L_{CD}}\,H_m\ }
\]
\end{keybox}

Esta relación la impone el \textbf{circuito}: es \textbf{lineal}, una recta por el origen
de pendiente negativa en el plano $(H, B)$ ---$H$ y $B$ del imán tienen signos
opuestos---. Se la llama recta \textbf{desmagnetizante}. Sola no alcanza: falta la
relación constitutiva del imán, su \textbf{curva de histéresis}, que es parte de
los datos (no es única: depende de la historia, y hay que decir en cuál se
está).

\begin{figure}[H]
  \centering
  % Clase 22 — solución gráfica del circuito con imán: intersección de la recta
% desmagnetizante B = -mu0 (L_AB/L_CD) H, impuesta por el circuito, con la curva
% de histéresis del imán (cualitativa). Hay dos soluciones simétricas porque no
% se dio la orientación del imán.
\begin{tikzpicture}
  \begin{axis}[
    emfig, width=9.5cm, height=7cm,
    axis lines=middle,
    xlabel style={anchor=north west, font=\small},
    ylabel style={anchor=south east, font=\small},
    xmin=-4.2, xmax=4.2, ymin=-1.65, ymax=1.65,
    xtick=\empty, ytick=\empty, grid=none, clip=false,
    xlabel={$H_m$}, ylabel={$B_m$},
  ]
    \addplot[figblue, domain=-4:4, samples=120] {1.3*tanh((x+1)/0.9)};
    \addplot[figblue, domain=-4:4, samples=120] {1.3*tanh((x-1)/0.9)};
    \addplot[figred, domain=-2.4:2.4, samples=2] {-0.6*x};
    % intersecciones (resueltas por biseccion: H = -0.6989, B = 0.4194)
    \fill[figred] (axis cs:-0.6989,0.4194) circle (2.2pt);
    \fill[figred] (axis cs:0.6989,-0.4194) circle (2.2pt);
    \node[font=\footnotesize, text=figred, anchor=north west] at (axis cs:2.45,-1.05)
      {$B_m = -\mu_0\dfrac{L_{AB}}{L_{CD}}\,H_m$};
  \end{axis}
\end{tikzpicture}

  \caption{Solución gráfica: la recta desmagnetizante que impone el circuito corta la curva de histéresis del imán (cualitativa) en dos puntos simétricos, porque no se dio la orientación del imán.}
\end{figure}

La solución está en la \textbf{intersección} de la recta con la curva de
histéresis. Hay \textbf{dos} soluciones, simétricas, porque no se dijo hacia dónde
apunta el imán (cuál es su polo norte).

\begin{notebox}[Cuándo no hay simetría]
Si además hay un bobinado, o dos imanes, ya
existe una dirección preferencial y las dos orientaciones no son
equivalentes. Hay que dar como dato hacia dónde apunta cada cosa; si no, hay
varias soluciones posibles. Un bobinado alrededor del imán no cambia nada:
lo que importa es cómo está enrollado.
\end{notebox}

En los ejercicios la curva de histéresis suele reemplazarse por curvas
poligonales para poder hacer las cuentas; si es una curva realista, hay que
trabajar con una tabla.

\subsection{Para qué sirven los circuitos magnéticos}

El uso más común son los \textbf{transformadores}: dos bobinados sobre el mismo
núcleo, el \textbf{primario}, por el que entra corriente, y el \textbf{secundario}, por
el que sale. El núcleo de hierro hace que todo el campo generado por uno llegue
al otro y no se pierda energía por esa vía. (Siempre hay pérdidas, pero las
principales son de naturaleza eléctrica, no magnética.)

Respondiendo a una pregunta: para el circuito magnético, un bobinado con
corriente impuesta desde afuera y uno cuya corriente es inducida entran igual
en las ecuaciones; en un caso la corriente es dato y en el otro, incógnita, que
se determina acoplando con el circuito eléctrico. El que tiene corriente
impuesta es el que aporta energía al sistema.

\section{Adelanto: inducción electromagnética}

Con los minutos que quedan, el docente abre el capítulo 8. En el caso estático
o estacionario,

\[
\curl\vec E = 0 \quad\Longleftrightarrow\quad \oint_C\vec E\cdot d\vec l = 0
\ \text{ para toda curva cerrada } C .
\]

Pero de Física III se sabe que, si el flujo magnético a través de $C$ varía con
el tiempo, vale la \textbf{ley de Faraday}:

\begin{keybox}
\[
\boxed{\ \oint_C\vec E\cdot d\vec l = -\frac{d\Phi}{dt}\ }
\]
\end{keybox}
con el flujo calculado sobre una superficie orientada con la regla de la mano
derecha respecto de $C$.

\begin{itemize}
  \item El caso estacionario es un caso particular: si el flujo no cambia, la
  circulación es cero.
  \item Si cambia, la fuerza eléctrica deja de ser conservativa: la integral de línea
  entre dos puntos depende del camino.
  \item Como siempre en este curso, se va a ir más allá de la forma integral, a una
  ecuación diferencial que dé el \textbf{rotacional} de $\vec E$ en términos de la
  variación del campo magnético.
\end{itemize}

\noindent\emph{La próxima clase retoma Faraday, con detalle sobre los signos y la regla de la
mano derecha.}

\end{document}
